% math4cs-hw1-prop-logic.tex

% !TEX program = xelatex
%%%%%%%%%%%%%%%%%%%%
% see http://mirrors.concertpass.com/tex-archive/macros/latex/contrib/tufte-latex/sample-handout.pdf
% for how to use tufte-handout
\documentclass[a4paper, justified]{tufte-handout}

\input{hw-preamble} % feel free to modify this file if you understand LaTeX well
%%%%%%%%%%%%%%%%%%%%
\title{计算机数学 (1-prop-logic: 命题逻辑)}
\me{魏恒峰}{hfwei@hnu.edu.cn}{}{}
\date{2026年3月12日}
%%%%%%%%%%%%%%%%%%%%
\begin{document}
\maketitle
%%%%%%%%%%%%%%%%%%%%
\noplagiarism % PLEASE DON'T DELETE THIS LINE!
%%%%%%%%%%%%%%%%%%%%
\begin{abstract}
  % \mfigcap{width = 0.85\textwidth}{figs/George-Boole}{George Boole}
  % \begin{center}{\fcolorbox{blue}{yellow!60}{\parbox{0.65\textwidth}{\large
  %   \begin{itemize}
  %     \item
  %   \end{itemize}}}}
  % \end{center}
\end{abstract}
%%%%%%%%%%%%%%%%%%%%
\beginrequired
%%%%%%%%%%%%%%%
\begin{problem}[命题逻辑公式上的数学归纳法]
  假设公式 $\alpha$ 中不含 ``$\lnot$'' 符号。
  请证明, $\alpha$ 中超过四分之一的符号是命题符号。

  ({\it 提示: 公式的长度定义见课件。})
\end{problem}

\begin{proof}
\end{proof}
%%%%%%%%%%%%%%%

%%%%%%%%%%%%%%%
\begin{problem}[命题逻辑公式的语义]
  请找出所有满足以下公式的真值指派:
  \begin{enumerate}[(1)]
    \item $\alpha: (A_{1} \to A_{2}) \land (A_{2} \to A_{3})
      \land \cdots \land (A_{n-1} \to A_{n})$
    \item $\beta: \alpha \land (A_{n} \to A_{1})$
  \end{enumerate}
\end{problem}

\begin{solution}

\end{solution}
%%%%%%%%%%%%%%%

%%%%%%%%%%%%%%%
\begin{problem}[重言蕴含]
  请利用重言式证明:
  \[
    \set{\alpha \to \beta} \models
      (\gamma \to \alpha) \to (\gamma \to \beta).
  \]
\end{problem}

\begin{solution}
\end{solution}
%%%%%%%%%%%%%%%

%%%%%%%%%%%%%%%
\begin{problem}[合取范式与析取范式]
  我们先引入一个定义。

  \begin{definition}[合取范式 (Conjunctive Normal Form; CNF)]
    \label{def:cnf}
    我们称公式 $\alpha$ 是\red{\bf 合取范式}, 如果它形如
    \[
      \alpha = \beta_{1} \land \beta_{2} \land \dots \land \beta_{k},
    \]
    其中, 每个 $\beta_{i}$ 都形如
    \[
      \beta_{i} = \beta_{i1} \lor \beta_{i2} \lor \dots \lor \beta_{in},
    \]
    并且 $\beta_{ij}$ 或是一个命题符号, 或是命题符号的否定。
  \end{definition}

  例如, 下面的公式就是一个合取范式。
  \[
    (P \lor \lnot Q \lor R) \;\red{\land}\; (\lnot P \lor Q) \;\red{\land}\; \lnot Q
  \]

  将定义~\ref{def:cnf} 中的所有$\land$换成$\lor$, 所有$\lor$换成$\land$, 其余不变,
  就变成了析取范式 (Disjunctive Normal Form; DNF)的定义。本题以CNF为例。

  将任意公式转化成CNF或DNF的方法如下:
  \begin{enumerate}[(1)]
    \item 先将公式中的联词化归成 $\lnot$, $\land$ 与 $\lor$;
    \item 再使用 De Morgan 律将 $\lnot$ 移到各个命题变元之前 (\blue{``否定深入''});
    \item 最后使用结合律、分配律将公式化归成合取范式或析取范式。
  \end{enumerate}

  请将
  \[
    (P \land (Q \to R)) \to S
  \]
  化为合取范式。
\end{problem}

\begin{solution}
\end{solution}
%%%%%%%%%%%%%%%

%%%%%%%%%%%%%%%
\begin{problem}[命题逻辑的自然推理系统]
  请使用自然推理规则证明以下结论:

  {提示: 你可以使用 \LaTeX{} package `logicproof' 书写漂亮的证明。}
  \begin{enumerate}[(1)]
    \item
      \[
        p  \to (q \to r) \vdash (p \land q) \to r
      \]
    \item
      \[
        (p \land q) \lor r \vdash (p \lor r) \land (q \lor r)
      \]
    \item
      \[
        (p \lor r) \land (q \lor r) \vdash (p \land q) \lor r
      \]
    \item
      \[
        \vdash \alpha \to (\beta \to \alpha)
      \]
    \item
      \[
        \vdash (\alpha \to (\beta \to \gamma)) \to ((\alpha \to \beta) \to (\alpha \to \gamma))
      \]
    \item
      \[
        \vdash (\lnot \beta \to \lnot \alpha) \to ((\lnot \beta \to \alpha) \to \beta)
      \]
  \end{enumerate}
\end{problem}

\begin{solution}

\end{solution}
%%%%%%%%%%%%%%%

%%%%%%%%%%%%%%%%%%%%
\beginoptional

%%%%%%%%%%%%%%%
\begin{problem}[命题逻辑的自然推理系统]
  请使用自然推理规则 (非 derived rules) 证明排中律:
  \[
    \vdash \alpha \lor \lnot\alpha
  \]
\end{problem}

\begin{solution}
\end{solution}
%%%%%%%%%%%%%%%

%%%%%%%%%%%%%%%
\begin{problem}[Build with AI]
  利用 AI 工具（如 Google AI Studio, Google Gemini,
  ChatGPT, Codex, Claude Code, Kimi 等）
  实现一个以``命题逻辑的自然推理系统''为主题的网页或应用程序。

  具体形式不限，可以是一个交互式的证明助手，
  也可以是证明通关游戏，等等。

  可以多人合作完成，提交时请注明参与者姓名。

  提交内容包括：
  \begin{itemize}
    \item 项目名称
    \item 项目作者
    \item 项目链接（如 GitHub 仓库链接，或在线应用链接）
    \item 项目截图（至少一张展示核心功能的截图）
    \item 项目简介（功能介绍，使用说明等）
    \item 与 AI 工具的核心对话记录
    \item 可选: 开发心得（开发过程中遇到的挑战、解决方案、感悟、AI 工具的表现等）
    \item 可选：项目演示视频链接（如 Bilibili 视频链接）
    \item 其它, 由你决定
  \end{itemize}
\end{problem}

\begin{solution}
\end{solution}
%%%%%%%%%%%%%%%%%%%%

%%%%%%%%%%%%%%%%%%%%
% 如果没有需要订正的题目，可以把这部分删掉

\begincorrection
%%%%%%%%%%%%%%%%%%%%

%%%%%%%%%%%%%%%%%%%%
% 如果没有反馈，可以把这部分删掉
\beginfb

你可以写 (也可以发邮件)
\begin{itemize}
  \item 对课程及教师的建议与意见
  \item 教材中不理解的内容
  \item 希望深入了解的内容
  \item $\cdots$
\end{itemize}
%%%%%%%%%%%%%%%%%%%%
\end{document}